\chapter{Quantile Regression}
\label{cap:qreg}
% ── index ───────────────────────────────────────────────────────
\index{quantile regression|textbf}
\index{qreg@\texttt{qreg()}|textbf}
\index{median}
\index{quantile}
\index{heteroscedasticity}

% ─────────────────────────────────────────────────────────────────────────────
\section{The model}

OLS estimates $E[y \mid x]$ --- the conditional mean. Quantile Regression
(QR) estimates $Q_\tau[y \mid x]$ --- the $\tau$-th conditional quantile for
any $\tau \in (0,1)$:

\[
  Q_\tau(y \mid x) = x'\beta(\tau)
\]

The vector $\beta(\tau)$ varies with $\tau$: each quantile has its own set
of coefficients. A single model can describe the entire conditional
distribution of $y$.

% ─────────────────────────────────────────────────────────────────────────────
\section{The estimator}

The estimator minimizes the \textbf{asymmetric loss function} (check function):

\[
  \hat\beta(\tau) = \arg\min_\beta \sum_{i=1}^{n} \rho_\tau(y_i - x_i'\beta)
\]

where
\[
  \rho_\tau(u) =
  \begin{cases}
    \tau\, u       & \text{if } u \geq 0 \\
    (\tau - 1)\,u  & \text{if } u < 0
  \end{cases}
\]

For $\tau = 0{,}5$: $\rho_{0,5}(u) = |u|/2$ --- minimization of the sum of
absolute deviations (LAD/median). For $\tau \neq 0{,}5$, positive and
negative residuals receive asymmetric weights, shifting the estimate toward the
$\tau$-th quantile.

\begin{nota}
  Unlike OLS, QR has no closed-form solution --- the problem
  is a linear program solved by the simplex method or interior-point methods.
  Therefore, standard inference uses bootstrap (\lstinline{boot=}).
\end{nota}

% ─────────────────────────────────────────────────────────────────────────────
\section{Advantages over OLS}

\begin{itemize}
  \item \textbf{Heteroscedasticity:} when $\mathrm{Var}[y \mid x]$ varies
        with $x$, the slopes $\beta_j(\tau)$ differ across quantiles. OLS
        captures only the implicit weighted average.

  \item \textbf{Robustness:} the check function is linear in the residuals, not
        quadratic. Outliers in $y$ influence $\hat\beta(\tau)$
        less than $\hat\beta_{\mathrm{OLS}}$.

  \item \textbf{Skewed distribution:} wages, income and asset prices
        have heavy right tails. QR describes effects in the tails without
        assumptions about the shape of the distribution.
\end{itemize}

% ─────────────────────────────────────────────────────────────────────────────
\section{DGP Verification}

The DGP below introduces multiplicative heteroscedasticity: the dispersion
of $y$ grows with $x$, so that the effect of $x$ is larger in upper quantiles
than in lower ones.

\[
  y_i = 2{,}0 + 1{,}5\,x_i + (1{,}0 + 0{,}8\,x_i)\,\varepsilon_i,
  \quad \varepsilon_i \sim \mathcal{N}(0,1)
\]

The conditional quantile $Q_\tau(y \mid x)$ is increasing in $x$, and the
slope $\partial Q_\tau / \partial x$ increases with $\tau$: individuals
in upper quantiles are more sensitive to $x$ than those in the lower quantile.

\begin{lstlisting}[caption={Heteroscedastic DGP --- OLS vs.\ four quantiles}]
seed(42)
input df
id
1
end
for i in 1..=10 { df = append(df, df) }
df |> filter(_n <= 1000)
generate df x = rnormal()
generate df e = rnormal()
generate df y = 2.0 + 1.5*x + (1.0 + 0.8*x)*e
let m_ols = ols(y ~ x, df)
let m_q25 = qreg(y ~ x, df, tau=0.25)
let m_q50 = qreg(y ~ x, df, tau=0.50)
let m_q75 = qreg(y ~ x, df, tau=0.75)
esttab(m_ols, m_q25, m_q50, m_q75)
\end{lstlisting}

\begin{lstlisting}[language={}, basicstyle=\ttfamily\small,
  frame=none, numbers=none, backgroundcolor=\color{codbg},
  xleftmargin=0pt, xrightmargin=0pt, breaklines=false]
                            (1)              (2)              (3)              (4)
                            OLS     QReg(τ=0.25)     QReg(τ=0.50)     QReg(τ=0.75)
──────────────────────────────────────────────────────────────────────────────────
x                     1.9947***        1.8739***        2.0149***        2.2462***
                       (0.0454)         (0.0687)         (0.0833)         (0.0579)
const                                  2.1659***        2.4525***        2.6637***
                                        (0.0302)         (0.0433)         (0.0369)
Constant              2.4538***
                       (0.0261)
──────────────────────────────────────────────────────────────────────────────────
N                          1000                0                0                0
R²                       0.6594
Adj. R²                  0.6591
──────────────────────────────────────────────────────────────────────────────────
* p<0.10  ** p<0.05  *** p<0.01
\end{lstlisting}

The slope of $x$ grows monotonically: $1{,}87$ at Q25, $2{,}01$ at Q50
and $2{,}25$ at Q75. OLS ($1{,}99$) is an implicit weighted average of these
heterogeneous effects --- it summarizes in a single number what QR decomposes
across the entire conditional distribution.

The intercept also increases with $\tau$ ($2{,}17 \to 2{,}45 \to 2{,}66$),
reflecting that upper quantiles of $y$ are above the mean for
any value of $x$.

% ─────────────────────────────────────────────────────────────────────────────
\section{Syntax}

\begin{lstlisting}
qreg(y ~ x1 + x2, df, tau=0.5)
qreg(y ~ x1 + x2, df, tau=0.5, boot=500)
\end{lstlisting}

The \lstinline{tau} parameter defines the desired quantile ($\tau \in (0,1)$,
default $0{,}5$). The \lstinline{boot} parameter activates bootstrap for
standard errors (default 200 replications).

\begin{lstlisting}[caption={Typical workflow --- wages at multiple quantiles}]
load "pnad.csv" as df

let q10 = qreg(wage ~ educ + exp, df, tau=0.10, boot=500)
let q50 = qreg(wage ~ educ + exp, df, tau=0.50, boot=500)
let q90 = qreg(wage ~ educ + exp, df, tau=0.90, boot=500)
let ols = ols(wage ~ educ + exp, df)

esttab(ols, q10, q50, q90)
\end{lstlisting}

\section{Capturing the result}

\begin{lstlisting}
let m = qreg(y ~ x, df, tau=0.75)
display m
let yhat = predict(m, "xb")
\end{lstlisting}
